\section*{Chapter \ref{ch:ml:low-latency-inference} --- Low-Latency Inference}

\begin{solution}{exo:ml:low-latency-inference:1}
Rounding to the nearest multiple of $s$ errs by at most $s/2 = \max|w|/254$: 0.39\% of the largest weight, and a far larger
share of a small weight.
\end{solution}

\begin{solution}{exo:ml:low-latency-inference:2}
$5\,738/544\,349 = 0.011$: the compiled forest's 99th percentile is about a ninety-fifth of the library call's.
\end{solution}

\begin{solution}{exo:ml:low-latency-inference:3}
A forest's prediction is a sum of leaf values chosen by comparisons: the compiled version makes the same comparisons and adds the
same doubles in the same order. A network's float arithmetic depends on the order of the sums in each product, on fused
multiply-adds and on vectorisation, which differ between libraries and compilers; parity is then a tolerance, unless the
arithmetic is integer, as in the int8 network.
\end{solution}

\begin{solution}{exo:ml:low-latency-inference:4}
$0.0123 = 0.7872\times2^{-6}$; $M = \mathrm{round}(0.7872\times2^{31}) = 1\,690\,499\,128$ and a total shift of $6 + 31 = 37$: the
product is $\mathrm{acc}\times M$ shifted right by 37 bits. Integers make the result identical on every platform and avoid
floating point on hardware that lacks it or where it is slower.
\end{solution}

\begin{solution}{exo:ml:low-latency-inference:5}
At eight bits the rounding is too fine to hurt, so \omterm{def:ml:representation-learning:ssl}{fine-tuning} only moves the model away from its trained optimum; at four bits
the rounding is coarse enough to damage the post-training model, and a model trained to expect it recovers more than half of the
loss (0.164 against 0.143, float 0.182).
\end{solution}

\begin{solution}{exo:ml:low-latency-inference:6}
Dividing a batch's time by its size gives the throughput's inverse, not the latency of a request: it hides the first request's
wait, the tail, and anything that happens once per call. Time single calls, as they arrive, and report the distribution.
\end{solution}

\begin{solution}{exo:ml:low-latency-inference:7}
Per-output scales give each row of weights its own grid, so small rows are not flattened by one large weight elsewhere: at four
bits the information coefficient rises from 0.143 to 0.153 (float 0.182), about a third of the loss recovered without any
retraining. At eight bits it is 0.182, as before.
\end{solution}

\begin{solution}{exo:ml:low-latency-inference:8}
With $x_i = s_xX_i$ and $w_i = s_wW_i$, $\sum_iw_ix_i + b = s_xs_w\big(\sum_iW_iX_i + b/(s_xs_w)\big)$: storing $B =
\mathrm{round}(b/(s_xs_w))$ in the accumulator's scale lets the bias be added to the integer sum. The output's integer is
$Y = \mathrm{round}\big((s_xs_w/s_y)(\sum_iW_iX_i + B)\big)$: the multiplier is $s_xs_w/s_y$, applied as a fixed-point product
and a shift.
\end{solution}

\begin{solution}{pb:ml:low-latency-inference:1}
\textbf{Part I.}
\begin{enumerate}
\item A 300-tree LightGBM forest (IC 0.224) and a $16\to32\to16\to1$ network (IC 0.182).
\item As flat node arrays (walked by the same loop in Python, C++ and Rust) and as generated C++ branches.
\item Trees: the same comparisons and the same order of double additions; the network: integer arithmetic.
\item The forest's 4\,200 splits and 4\,500 leaves in arrays of a few hundred kilobytes of text and far less in binary; the
network 1\,236 bytes in int8 against 4\,356 in float32.
\end{enumerate}
\textbf{Part II.}
\begin{enumerate}[resume]
\item Median and 99th percentile (µs): LightGBM 265 and 544; Python forest 300 and 397; PyTorch network 21 and 44; NumPy int8
network 24 and 66; C++ forest loop 6.2 and 13.2; C++ branches 2.6 and 5.7; C++ int8 network 0.47 and 1.1; Rust forest 7.8
and 17.8; Rust int8 network 0.58 and 1.9.
\item The interpreter, argument checks, conversions of the row to and from arrays, and thread-pool machinery; the trees
themselves are a few microseconds.
\item Thresholds and branch targets are in the code, not loaded from memory, and the compiler lays out the paths.
\item About a quarter, for bounds checks on every array access; removable after validating the model at load.
\end{enumerate}
\textbf{Part III.}
\begin{enumerate}[resume]
\item Nothing measurable after training (0.182); quantisation-aware \omterm{def:ml:representation-learning:ssl}{fine-tuning} 0.179.
\item \omterm{def:ml:low-latency-inference:quant}{Post-training quantisation} falls to 0.143, \omterm{def:ml:low-latency-inference:quant}{quantisation-aware training} holds 0.164, per-channel scales 0.153.
\item From 273~µs per call for one row to 3.1~ms for 1\,024: 3.0~µs per row.
\item For scoring many requests with no latency constraint (a universe each second); not for a decision on one instrument.
\end{enumerate}
\textbf{Part IV.}
\begin{enumerate}[resume]
\item \emph{Two microseconds for a forest.} The generated C++ forest's 99th percentile is 5.7~µs against 544~µs for the library
call, a ratio of about 1 to 95 (2.6 against 265~µs at the median); int8 \omterm{def:ml:low-latency-inference:quant}{quantisation} costs the network nothing measurable after
training (0.182) and 0.003 with quantisation-aware \omterm{def:ml:representation-learning:ssl}{fine-tuning} (0.179); at four bits, 0.040 and 0.018.
\item The forest as generated C++ (2.6~µs median, 5.7 at the tail) if its extra signal is worth it, or the int8 network at half a
microsecond.
\item The source model, the compiled artefact and its hash, the compiler and flags, and the parity report.
\item Isolated cores, no other jobs, warm caches as in production, single requests at the production rate, and tails over long
runs, with the measurement's own overhead measured.
\item Ten times the code: the generated source may no longer fit the instruction cache, and flat arrays or vectorised evaluation
may win.
\item Where the budget is below a microsecond and fixed, and the model is small and stable enough to be worth a hardware cycle.
\item As a difference between the features the compiled model receives and those the research model saw, invisible to the
kernel's parity tests: parity must also be checked on live inputs.
\item Mostly around the model, not in it.
\end{enumerate}
\end{solution}

\begin{solution}{iq:ml:low-latency-inference:1}
Export the trees and compile them (generated branches or flat arrays in C++), keep the process hot, avoid allocation and
conversions on the path, check parity against the library, and measure tails one request at a time.
\iqlookfor{compilation, a clean hot path, parity and tail measurement.}
\end{solution}

\begin{solution}{iq:ml:low-latency-inference:2}
Real values are integers times a scale (plus a zero point in asymmetric schemes); products accumulate in int32; each layer's
accumulator is rescaled to the next layer's integers by a fixed-point multiplier; biases are stored at the accumulator's scale.
\iqlookfor{scales, accumulators and requantisation.}
\end{solution}

\begin{solution}{iq:ml:low-latency-inference:3}
No heap allocation, locks, system calls, exceptions or virtual dispatch; bounded loops; data laid out for the cache; everything
loaded and validated at start-up.
\iqlookfor{the usual hot-path rules with reasons.}
\end{solution}

\begin{solution}{iq:ml:low-latency-inference:4}
Shared test vectors with reference outputs from the training framework; exact comparison for trees and integer networks, stated
tolerances for float ones; the compiled artefact's hash registered with the source model; checks on live inputs too.
\iqlookfor{shared vectors, exactness where possible, registration.}
\end{solution}

\begin{solution}{iq:ml:low-latency-inference:5}
It helps when per-call overhead dominates and requests can wait; it hurts when each request has a deadline, because the first
request in a batch waits for the last.
\iqlookfor{overhead amortisation against waiting time.}
\end{solution}

\begin{solution}{iq:ml:low-latency-inference:6}
Quantise after training and measure the metric that matters; if the loss is negligible, stop; if not, try finer granularity
(per channel) and, if needed, quantisation-aware \omterm{def:ml:representation-learning:ssl}{fine-tuning}, and compare on held-out data.
\iqlookfor{measure first, then escalate.}
\end{solution}
